% Streng antichronologisch. Nespresso laeuft noch und steht deshalb oben, auch
% wenn PwC der staerkere Eintrag ist. Die Reihenfolge zu brechen, um sich besser
% darzustellen, faellt einem deutschen Leser sofort auf.
%
% "Eigenstaendige Durchfuehrung" ist gestrichen. Ein Praktikant und Werkstudent
% fuehrt keine Jahresabschlusspruefung eigenstaendig durch, und wer die Branche
% kennt, rechnet solche Formulierungen heraus.
\section*{Berufserfahrung}
\resumeSubHeadingListStart

\resumeSubheading
{Nespresso (Nestlé)}{Köln / Düsseldorf}
{Produktberater (Teilzeit)}{seit Jan.\@ 2026}
\resumeItemListStart
\resumeItem{Beratung und Verkauf von Premium-Produkten im direkten Kundenkontakt, bei konsistenter Zielerreichung.}
\resumeItemListEnd

\resumeSubheading
{PwC Deutschland}{Düsseldorf, Deutschland}
{Praktikant \& Werkstudent, Assurance Solutions}{März 2023 -- Juli 2025}
\resumeItemListStart
\resumeItem{Abweichungsanalysen und Plausibilitätsprüfungen bei Jahresabschlussprüfungen (HGB/IFRS) für Mandate mit Bilanzsummen von über 500\,Mio.\,EUR.}
\resumeItem{Auswertung und Aufbereitung großer Datenbestände in Excel, darunter KPI-Tracking und Forecasting für das Prüfungsteam.}
\resumeItem{Identifikation wesentlicher Geschäftstreiber und Risiken gemeinsam mit Senior Associates und Managern. Die Ergebnisse gingen direkt in die Prüfungsberichte ein.}
\resumeItem{Mitarbeit an Prüfungsberichten und Management Letters für Mandanten.}
\resumeItemListEnd

\resumeSubHeadingListEnd
